El radó és un gas noble radioactiu que es produeix de manera natural al subsol com a resultat de la cadena de desintegració de l'urani. El radó es pot acumular en espais tancats com ara soterranis. La inhalació de concentracions superiors a $300\ \mathrm{Bq\,m^{-3}}$ d'aquest gas durant períodes prolongats incrementa significativament el risc de patir càncer de pulmó.

\begin{apartat}{1,25}
Observeu la figura i escriviu les principals desintegracions que constitueixen la cadena de desintegració de l'urani fins a la formació del radó. En cada cas, indiqueu de quin tipus de reacció es tracta.
\end{apartat}

\begin{apartat}{1,25}
Es prenen mesures de la massa i l'activitat del radó en una habitació tancada de $18\ \mathrm{m^3}$ de volum. Si es mesura una massa de 0,80~pg de radó, a quina activitat correspon? Aquest valor està per sobre o per sota del llindar d'activitat per unitat de volum indicat a l'enunciat? Si inicialment es mesura una activitat per unitat de volum de $230\ \mathrm{Bq\,m^{-3}}$ i l'habitació queda perfectament aïllada, quina activitat per unitat de volum hi haurà al cap de 15 dies?
\end{apartat}

\figura[0.5]{fig1.png}

\dades{%
  Nombre d'Avogadro: $N_\mathrm{A} = 6,022\times 10^{23}$.\\
  $m({}^{222}_{\ 86}\mathrm{Rn}) = 222\ \mathrm{g\,mol^{-1}}$.\\
  El període de semidesintegració del ${}^{222}_{\ 86}\mathrm{Rn}$ és de 3,82 dies.}
